\documentclass[10pt, twocolumn, letterpaper]{article}

% --- Essential Packages ---
\usepackage[utf8]{inputenc}
\usepackage[T1]{fontenc}
\usepackage{libertine}
\usepackage[libertine]{newtxmath}
\usepackage{microtype}
\usepackage[hmargin=1.8cm, vmargin=2.0cm, columnsep=0.6cm]{geometry}
\usepackage{titlesec}
\usepackage{booktabs}
\usepackage{graphicx}
\usepackage{xcolor}
\usepackage{hyperref}
\usepackage{enumitem}
\usepackage{caption}
\usepackage{amsmath}
\usepackage{float} % Added to strictly control figure placement
\usepackage{colortbl} % Row shading in tables

% --- Aesthetic Customization ---
\definecolor{DeepTeal}{HTML}{004D40}

\hypersetup{
    colorlinks=true,
    linkcolor=DeepTeal,
    citecolor=DeepTeal,
    urlcolor=DeepTeal
}

\titleformat{\section}
  {\color{DeepTeal}\normalfont\large\bfseries\scshape}{\thesection}{1em}{}[{\titlerule[0.5pt]}]
\titleformat{\subsection}
  {\color{DeepTeal}\normalfont\normalsize\bfseries}{\thesubsection}{1em}{}

\captionsetup{font=small, labelfont=bf, skip=4pt}
\setlength{\abovecaptionskip}{4pt}
\setlength{\belowcaptionskip}{0pt}
\raggedbottom % Prevents awkward vertical stretching across columns

% --- Document Info ---
\title{
    \vspace{-1.5em}
    \rule{\linewidth}{1.5pt} \\
    \Large \textbf{CS-233: Introduction to Machine Learning} \\
    \normalsize \textit{Project Milestone 1 Report} \\
    \vspace{-0.5em}
    \rule{\linewidth}{1.5pt}
}

\author{
    \begin{tabular}{c c c}
        \textbf{Author 1} & \textbf{Author 2} & \textbf{Author 3} \\
        SCIPER: 123456 & SCIPER: 234567 & SCIPER: 345678 \\
        \textit{name1@epfl.ch} & \textit{name2@epfl.ch} & \textit{name3@epfl.ch}
    \end{tabular}
}
\date{}

\begin{document}

\maketitle

\section{Introduction}
We evaluate traditional machine learning algorithms on the Gaming and Mental Health Dataset (2000 samples, 13 features). Milestone 1 addresses two supervised tasks:
\begin{itemize}[noitemsep, leftmargin=*]
    \item \textbf{Regression:} Predicting a continuous addiction score (0--10).
    \item \textbf{Classification:} Predicting discrete addiction levels (Low, Medium, High).
\end{itemize}
We implemented Linear Regression, multiclass Logistic Regression, and K-Nearest Neighbors (KNN) from scratch.

\section{Method}
\subsection{Data Preparation}
All models share a unified preprocessing pipeline:
\begin{itemize}[noitemsep, leftmargin=*]
    \item \textbf{Validation Split:} 70\% training / 30\% validation. Final evaluations run on a separate 400-sample test set.
    \item \textbf{Normalization:} Features are standardized (zero mean, unit variance) using training set statistics.
    \item \textbf{Bias Augmentation:} A column of ones is prepended for Linear and Logistic Regression.
\end{itemize}

\subsection{Algorithms}
\begin{itemize}[noitemsep, leftmargin=*]
    \item \textbf{Linear Regression:} Computed via the closed-form pseudo-inverse solution ($\mathbf{w} = \mathbf{X}^{\dagger}\mathbf{y}$).
    \item \textbf{Logistic Regression:} Multiclass softmax classifier trained by full-batch gradient descent. A row-max shift is applied to logits to prevent numerical overflow.
    \item \textbf{KNN:} Relies on Euclidean distance. Classification uses a majority vote among $k$ neighbors; regression averages their target values.
\end{itemize}

% --- Force break to start the next section at the top of the second column ---
\newpage

\section{Experiments and Results}
Models were evaluated using Mean Squared Error (MSE) for regression, and Accuracy alongside Macro-F1 for classification to account for class imbalances.

\subsection{Hyperparameter Search}
We tuned hyperparameters on the validation split.

\vspace{0.5em}
\noindent \textbf{KNN ($k$):} Tested $k \in \{1, 3, 5, 7, 9, 11, 15, 21\}$. As seen in Figure~\ref{fig:knn}, $k=9$ is the optimal parameter, maximizing Classification Accuracy (80.6\%) and Macro-F1, while minimizing Regression MSE (1.947). Smaller $k$ severely overfit, while larger $k$ over-smooth.

\begin{figure}[H]
\centering
\includegraphics[width=\linewidth]{figures/knn_tuning.png}
\caption{KNN validation curves. \textbf{Left:} Classification metrics peak at $k{=}9$. \textbf{Right:} Regression MSE is minimized at $k{=}9$.}
\label{fig:knn}
\end{figure}

\noindent \textbf{Logistic Regression ($\eta$, iters):} The optimal configuration is $\eta=0.3$ with 500 iterations (Figure~\ref{fig:logreg}), achieving 86.04\% validation accuracy. Lower learning rates systematically underfit within the budget.

\begin{figure}[H]
\centering
\includegraphics[width=\linewidth]{figures/logreg_heatmap.png}
\caption{LogReg validation search. \textbf{Left:} Accuracy (\%). \textbf{Right:} Macro-F1. Both metrics are best at $\eta{=}0.3$, 500 iterations (red border).}
\label{fig:logreg}
\end{figure}

% --- Force break to push the next subsection to the new page ---
\newpage

\subsection{Final Test Performance}
The optimal models were retrained and evaluated on the held-out test set. Table~\ref{tab:results} summarizes the final metrics. Figure~\ref{fig:cm} and Figure~\ref{fig:dist} illustrate the classification behaviors across the imbalanced dataset.

\begin{table}[H]
\centering
\caption{Final method comparison on validation and test sets. Shaded rows are the best method per metric.}
\label{tab:results}
\setlength{\tabcolsep}{5pt}
\renewcommand{\arraystretch}{1.15}
\begin{tabular}{@{}llcc@{}}
\toprule
\textbf{Method} & \textbf{Metric} & \textbf{Val.} & \textbf{Test} \\
\midrule
\multicolumn{4}{@{}l}{\textit{Regression}} \\
\rowcolor{DeepTeal!15}
Linear Regression        & MSE $\downarrow$ & 0.931 & 0.994 \\
KNN Reg.\ ($k{=}9$)     & MSE $\downarrow$ & 1.947 & 1.865 \\
\midrule
\multicolumn{4}{@{}l}{\textit{Classification}} \\
\rowcolor{DeepTeal!15}
LogReg ($\eta{=}0.3$, 500) & Acc.\,(\%) $\uparrow$ & 86.04 & 88.25 \\
KNN Cls.\ ($k{=}9$)      & Acc.\,(\%) $\uparrow$ & 80.63 & 79.75 \\
\rowcolor{DeepTeal!15}
LogReg ($\eta{=}0.3$, 500) & Macro-F1 $\uparrow$  & 0.727 & 0.819 \\
KNN Cls.\ ($k{=}9$)      & Macro-F1 $\uparrow$  & 0.579 & 0.590 \\
\bottomrule
\end{tabular}
\end{table}

\begin{figure}[H]
\centering
\includegraphics[width=0.88\linewidth]{figures/confusion_matrix.png}
\caption{Confusion matrix for LogReg ($\eta{=}0.3$, 500 iters) on the test set. Two key errors: \emph{Medium}$\to$\emph{Low} (29/107, frequent due to class imbalance) and \emph{High}$\to$\emph{Medium} (5/13, important as severe cases are downgraded).}
\label{fig:cm}
\end{figure}

\begin{figure}[H]
\centering
\includegraphics[width=\linewidth]{figures/class_distribution.png}
\caption{Class distribution showing a heavy skew toward the \emph{Low} addiction class in both splits.}
\label{fig:dist}
\end{figure}
\newpage

\subsection{Runtime Analysis}
Training and prediction times were measured on the full 1600-sample training set and 400 test samples (13 features).

\begin{table}[H]
\centering
\caption{Average training and prediction times.}
\label{tab:runtime}
\setlength{\tabcolsep}{4pt}
\begin{tabular}{@{}lcc@{}}
\toprule
\textbf{Method} & \textbf{Train} & \textbf{Predict (400)} \\
\midrule
Linear Regression          & 0.43\,ms  & 0.01\,ms \\
Logistic Reg.\ (500 iter.) & 119\,ms   & 0.09\,ms \\
KNN ($k{=}9$)              & $\approx$0 & $\sim$35\,ms \\
\bottomrule
\end{tabular}
\end{table}

\begin{itemize}[noitemsep, leftmargin=*]
    \item \textbf{Linear Regression:} Fastest overall — trains in under 1\,ms with equally instant prediction, as no iterative optimization is needed.
    \item \textbf{Logistic Regression:} Takes $\sim$119\,ms to train (500 gradient steps on 1600 samples), but predicts in under 0.1\,ms once the weights are learned.
    \item \textbf{KNN:} Requires no training at all (it simply stores the data), but is the slowest at prediction ($\sim$35\,ms) since it must scan all training points for every new query.
\end{itemize}

\section{Discussion and Conclusion}
\begin{itemize}[leftmargin=*]
    \item \textbf{Regression:} Linear Regression significantly outperformed KNN (Test MSE 0.994 vs.\ 1.865), confirming a strong underlying linear relationship in the features. KNN is also penalized by having 13 features: as dimensionality grows, all points tend to become roughly equidistant, so the nearest neighbors are not necessarily the most similar ones, weakening the voting mechanism.
    \item \textbf{Classification:} Logistic Regression is vastly superior to KNN (88.25\% accuracy, F1 0.819 vs.\ 0.590). This is expected: Logistic Regression learns a global discriminative boundary from all samples, while KNN makes local distance-based decisions that are less reliable in 13-dimensional, imbalanced data where minority-class neighbors are sparse. From Figure~\ref{fig:cm}, two errors stand out: \emph{Medium}$\to$\emph{Low} (29/107, driven by class imbalance) and \emph{High}$\to$\emph{Medium} (5/13), where severe cases are missed --- their low recall ($8/13 \approx 62\%$) reduces Macro-F1, which is why it is reported alongside accuracy.
\end{itemize}

\paragraph{Conclusion:} Linear Regression and Logistic Regression are the most robust methods for their respective tasks on this dataset. Future work for Milestone 2 will explore regularization techniques and feature engineering to further improve boundary classifications.

\end{document}